%+PATCH,//BRILL/TEX
%+DECK,abbrevs,T=LATEX.
\newcommand{\vitemt}[2]{
\begin{center}
{
\setlength{\fboxrule}{3pt}
\fcolorbox{schwachgrau}{hellgrau}
  {\color{black}
    \setlength{\fboxrule}{0pt}
    \framebox[0.5\textwidth]{#1 \hfill}
    \colorbox{white}{\color{black} #2}
  }
}
\end{center}
}

\newcommand{\vitem}[2]{
\begin{center}
{
\setlength{\fboxrule}{3pt}
\fcolorbox{schwachgrau}{hellgrau}
  {\color{black}
    \setlength{\fboxrule}{0pt}
    \framebox[0.25\textwidth]{\small #1 \hfill}
    \framebox[0.2\textwidth]{\hfill
      \colorbox{white}{\color{black} \small #2}
    }
  }
}
\end{center}
}

\newcommand{\bsp}{\hspace{-0.2cm}}

\newcommand{\git}[1]{\textbf{\color{grau}\textit{#1}} \color{black}}
\newcommand{\gits}[1]{\textbf{\color{grau}\textit{\textsmaller[0.3]{#1}}} \color{black}}
\newcommand{\gitse}[1]{\textbf{\color{grau}\textit{\textsmaller[0.3]{#1}}}\color{black}}
\newcommand{\gite}[1]{\textbf{\color{grau}\textit{#1}}\color{black}}

\newcommand{\gti}[1]{\textbf{\color{grau}\textit{#1}} \color{black}}
\newcommand{\gtis}[1]{\textbf{\color{grau}\textit{\textsmaller[0.3]{#1}}} \color{black}}
\newcommand{\gtise}[1]{\textbf{\color{grau}\textit{\textsmaller[0.3]{#1}}}\color{black}}
\newcommand{\gtie}[1]{\textbf{\color{grau}\textit{#1}}\color{black}}

\newcommand{\witem}[1]{
  \begin{center}
    \colorbox{schwachgrau}{
      \colorbox{hellgrau}{
      \color{grau}
      \framebox[0.45\textwidth]{\color{black} \small #1}
      }
    }
  \end{center}
}

\newcommand{\witemrot}[1]{
  \begin{center}
    \colorbox{yellow}{
      \colorbox{red}{
      \color{red}
      \framebox[0.45\textwidth]{\color{white} \small #1}
      }
    }
  \end{center}
}

\newcommand{\witems}[1]{
  \begin{center}
    \colorbox{yellow}{
      \colorbox{blue}{
      \color{blue}
      \framebox[0.5\textwidth]{\color{white} #1}
      }
    }
  \end{center}
}

\newcommand{\witemss}[1]{
  \begin{center}
    \colorbox{yellow}{
      \colorbox{blue}{
      \color{blue}
      \framebox[0.5\textwidth]{\color{white} #1}
      }
    }
  \end{center}
}

\newcommand{\witemsss}[1]{
  \begin{center}
    \colorbox{yellow}{
      \colorbox{blue}{
      \color{blue}
      \framebox[0.5\textwidth]{\color{white} #1}
      }
    }
  \end{center}
}

\newcommand{\ryellow}[1]{
  \begin{center}
    \colorbox{yellow}{\color{yellow}
      \framebox[0.6\textwidth]{\color{black} \small #1}
    }
  \end{center}
}
\newcommand{\rorange}[1]{
  \begin{center}
    \colorbox{orange}{\color{orange}
      \framebox[0.6\textwidth]{\color{black} \small #1}
    }
  \end{center}
}
\newcommand{\ryellows}[1]{
  \begin{center}
    \colorbox{yellow}{\color{yellow}
      \framebox[0.55\textwidth]{\color{black} \small #1}
    }
  \end{center}
}
\newcommand{\roranges}[1]{
  \begin{center}
    \colorbox{orange}{\color{orange}
      \framebox[0.55\textwidth]{\color{black} \small #1}
    }
  \end{center}
}
\newcommand{\rmft}[1]{\ryellow{magnetic field and trajectory}}
\newcommand{\rby}[1]{\ryellows{By(x)}}
\newcommand{\rbys}[1]{\ryellows{By(x) with observed sources of radiation}}
\newcommand{\rnz}[1]{\ryellows{nz(x)}}
\newcommand{\rnzs}[1]{\ryellows{nz(x) with observed sources of radiation}}
\newcommand{\rnzp}[1]{\ryellows{nzp(x)}}
\newcommand{\rnzps}[1]{\ryellows{nzp(x) with observed sources of radiation}}

\newcommand{\witemold}[1]{
  \begin{center}
    \colorbox{yellow}{
      \colorbox{blue}{
      \color{blue}
      \framebox[0.6\textwidth]{\color{white} \small #1}
      }
    }
  \end{center}
}

\newcommand{\witemsold}[1]{
  \begin{center}
    \colorbox{yellow}{
      \colorbox{blue}{
      \color{blue}
      \framebox[0.575\textwidth]{\color{white} \small #1}
      }
    }
  \end{center}
}

\newcommand{\witemssold}[1]{
  \begin{center}
    \colorbox{yellow}{
      \colorbox{blue}{
      \color{blue}
      \framebox[0.55\textwidth]{\color{white} \small #1}
      }
    }
  \end{center}
}

\newcommand{\witemsssold}[1]{
  \begin{center}
    \colorbox{yellow}{
      \colorbox{blue}{
      \color{blue}
      \framebox[0.50\textwidth]{\color{white} \small #1}
      }
    }
  \end{center}
}

\definecolor{darkslategray}{rgb}{0.18, 0.31, 0.31}
\definecolor{gray}{rgb}{0.18, 0.31, 0.31}

\definecolor{orange}{rgb}{1.,0.8,0.4} % rgb
\definecolor{hellgrau}{gray}{0.9}
\definecolor{schwachgrau}{gray}{0.95}
\definecolor{grau}{gray}{0.4}
\definecolor{dunkelgrau}{gray}{0.2}

\newcommand{\vitemold}[2]{
\begin{center}
  \colorbox{yellow}{
    \color{yellow}
    \colorbox{white}{
      \color{white}
      \framebox[0.6\textwidth]{
        \colorbox{blue}{\color{white} \small #1}
        \hfill \
        \colorbox{blue}{\color{white} \small #2}
      }
    }
  }
\end{center}
}

\newcommand{\grau}{\color{grau}}
\newcommand{\dunkelgrau}{\color{dunkelgrau}}
\newcommand{\schwarz}{\color{black}}

\newcommand{\Lin}{{\textit{\grau{\textbf{Linux~}}}}}
\newcommand{\bash}{{\textit{\grau{\textbf{bash~}}}}}
\newcommand{\Win}{{\textit{\grau{\textbf{Windows~}}}}}
\newcommand{\Lie}{{\textit{\grau{\textbf{Linux}}}}}
\newcommand{\Wine}{{\textit{\grau{\textbf{Windows}}}}}

\newcommand{\PyB}{{\textit{\grau{\textbf{PyBrill~}}}}}
\newcommand{\pyB}{{\textit{\grau{\textbf{pyBrill~}}}}}
\newcommand{\pyBe}{{\textit{\grau{\textbf{pyBrill}}}}}

\newcommand{\PB}{{\textit{\grau{\textbf{PyBrill~}}}}}
\newcommand{\pB}{{\textit{\grau{\textbf{pyBrill~}}}}}
\newcommand{\pBe}{{\textit{\grau{\textbf{pyBrill}}}}}

\newcommand{\bri}{{\textit{\grau{\textbf{brill~}}}}}
\newcommand{\brie}{{\textit{\grau{\textbf{brill}}}}}

\newcommand{\sta}{{\textit{\grau{\textbf{stage~}}}}}
\newcommand{\stae}{{\textit{\grau{\textbf{stage}}}}}

\newcommand{\bin}{{\textit{\grau{\textbf{bin}}}}}
\newcommand{\bine}{{\textit{\grau{\textbf{bin}}}}}

\newcommand{\wdir}{{\textit{\grau{\textbf{working directory~}}}}}
\newcommand{\wdire}{{\textit{\grau{\textbf{working directory}}}}}

\newcommand{\Py}{{\textit{\grau{\textbf{Python~}}}}}
\newcommand{\Pye}{{\textit{\grau{\textbf{Python}}}}}
\newcommand{\F}{{\textit{\grau{\textbf{FORTRAN~}}}}}
\newcommand{\Fe}{{\textit{\grau{\textbf{FORTRAN}}}}}

\newcommand{\Urp}{{\textit{\grau{\textbf{Urad\_phase~}}}}}
\newcommand{\urp}{{\textit{\grau{\textbf{urad\_phase~}}}}}
\newcommand{\urpe}{{\textit{\grau{\textbf{urad\_phase}}}}}

\newcommand{\Unam}{{\textit{\grau{\textbf{Urad\_phase.nam~}}}}}
\newcommand{\unam}{{\textit{\grau{\textbf{urad\_phase.nam~}}}}}
\newcommand{\uname}{{\textit{\grau{\textbf{urad\_phase.nam}}}}}

\newcommand{\mhb}{{\textit{\grau \textsmaller[2]{\textbf{WAVE.mhb~}}}}}
\newcommand{\mhbe}{{\textit{\grau \textsmaller[2]{\textbf{WAVE.mhb}}}}}
\newcommand{\dw}{{\textit{\grau \textsmaller[1.]{\$WAVE}}}}
\newcommand{\WSH}{{\textit{\grau\textsmaller[2]{\textbf{WAVESHOP~}}}}}
\newcommand{\WSHE}{{\textit{\grau\textsmaller[2]{\textbf{WAVESHOP}}}}}

\newcommand{\scap}[1]{\textit{\textsmaller[1]{\textbf{\grau #1}}}}
\newcommand{\cbox}[1]{\vspace{0.1cm}\hspace{1cm}\parbox[c]{0.8\textwidth}{\gti{\hspace{-0.2cm}#1}}}
\newcommand{\cboxs}[1]{\hspace{1cm} \parbox[c]{0.8\textwidth}{\grau \small \gti{#1}}}
\newcommand{\cboxss}[1]{\hspace{1cm}\parbox[c]{0.8\textwidth}{\grau
\textbf{\textsmaller[2]{#1}}}}
\newcommand{\cboxsl}[1]{\parbox[c]{0.9\textwidth}{\small \gti{#1}}}
\newcommand{\hin}{\hspace{1cm}}
\newcommand{\ro}{{\textit{source \dw/shell/wave\_root.sh}}}
\newcommand{\row}{{\textit{source \dw/shell/wave\_root.sh wave}}}
\newcommand{\gui}{{\textit{source \dw/shell/wave\_root.sh gui}}}
\newcommand{\wi}{{\gtis{wave.in~}}}
\newcommand{\wie}{{\gtise{wave.in}}}
\newcommand{\wo}{{\gtis{wave.out~}}}
\newcommand{\woe}{{\gtis{wave.out}}}
\newcommand{\wvs}{{\gti{wave.wvs~}}}
\newcommand{\wpl}{{\gti{wplot~}}}
\newcommand{\wple}{{\gti{wplot \bsp}}}
\newcommand{\wpg}{{\gti{waveplot.py~}}}
\newcommand{\wro}{{\gti{WAVE.root~}}}
\newcommand{\w}{{\gti{wave}}}
\newcommand{\senv}{\gtis{. \dw/shell/set\_wave\_environment.sh}}
\newcommand{\we}{{\gtis{\dw/input/wave.examples~}}}
\newcommand{\wee}{{\gtis{\dw/input/wave.examples}}}
\newcommand{\U}{{\small UNDUMAG~}}
\newcommand{\Ue}{{\small UNDUMAG}}
\newcommand{\W}{{\small WAVE~}}
\newcommand{\We}{{\small WAVE}}
\newcommand{\WS}{{{\small WAVES~}}}
\newcommand{\R}{{\small ROOT}}
\newcommand{\ws}{{\textit{waves~}}}
%\newcommand{\wpl}{{\textit{waves.pl~}}}
%\newcommand{\wpy}{{\textit{waves.py~}}}
\newcommand{\ispecnull}{{{\scap{ISPEC=0~}}}}
\newcommand{\ispecone}{{{\scap{ISPEC=1~}}}}
\newcommand{\ispec}{{\grau \bf \scap{ISPEC~}}}
\newcommand{\imagsplnnull}{{\scap{IMAGSPLN=0~}}}
\newcommand{\imagsplnone}{{\scap{IMAGSPLN=1~}}}
\newcommand{\imagspln}{{\scap{IMAGSPLN~}}}
\newcommand{\ieneloss}{{\scap{IENELOSS~}}}
\newcommand{\xstart}{{\scap{XSTART~}}}
\newcommand{\ystart}{{\scap{YSTART~}}}
\newcommand{\zstart}{{\scap{ZSTART~}}}
\newcommand{\vxin}{{\scap{VXIN~}}}
\newcommand{\vyin}{{\scap{VYIN~}}}
\newcommand{\vzin}{{\scap{VZIN~}}}
\newcommand{\dmyenergy}{{\scap{DMYENERGY~}}}
\newcommand{\dmycurr}{{\scap{DMYCURR~}}}
\newcommand{\ienelossnull}{{\scap{IENELOSS=0~}}}
\newcommand{\ienelossone}{{\scap{IENELOSS=1~}}}
\newcommand{\dend}{{\textbf{\grau\bf\textit{\textsmaller[1]{\$END~}}}}\schwarz}
\newcommand{\dende}{{\textbf{\grau\bf\textit{\textsmaller[1]{\$END}}}}\schwarz}
\newcommand{\contrl}{{\textbf{\grau\bf\textit{\textsmaller[1]{\$CONTRL~}}}}\schwarz}
\newcommand{\contrle}{{\textbf{\grau\bf\textit{\textsmaller[1]{\$CONTRL}}}}\schwarz}
\newcommand{\halbach}{{{\scap{\$HALBACH~}}}}
\newcommand{\ellip}{{{\scap{\$ELLIP~}}}}
\newcommand{\halbasy}{{{\scap{\$HALBAS}Y~}}}

%\renewcommand{\chaptername}{}
%\renewcommand{\contentsname}{\vspace{-1cm} \center{\Large Contents} \\[-1cm]}
%\renewcommand{\bibname}{\vspace{-3cm} \Large Literature \\[-0.4cm]}
\renewcommand{\listfigurename}{\vspace{-3cm} \Large Figures \\[-0.4cm]}
\renewcommand{\listtablename}{\vspace{-3cm} \Large Tables \\[-0.4cm]}
\renewcommand{\appendixname}{Appendix}
\renewcommand{\figurename}{Fig.}
\renewcommand{\tablename}{Tab.}

\newcommand{\lchib}{\mbox{ $ \large{\bar{ \chi}} $ }}
\newcommand{\chib}{\mbox{ $ \bar{ \chi } $ }}

\newcommand{\bt}{\tilde{\beta}}
\newcommand{\bnt}{\tilde{\beta_0}}
\newcommand{\boptt}{\tilde{\beta_{opt}}}
\newcommand{\bnopt}{{\beta}_{0,opt}}
\newcommand{\bnoptt}{\tilde{\beta}_{0,opt}}
\newcommand{\et}{\tilde{\eta}}
\newcommand{\eet}{\tilde{\eta_e}}
\newcommand{\eeopt}{\eta_{e,opt}}
\newcommand{\eeoptt}{\tilde{\eta}_{e,opt}}
\newcommand{\eit}{\tilde{\eta_i}}
\newcommand{\AT}{\tilde{A}}
\newcommand{\BT}{\tilde{B}}
\newcommand{\CT}{\tilde{C}}
\newcommand{\DT}{\tilde{D}}
\newcommand{\DOPTT}{\tilde{D}_{opt}}
\newcommand{\ET}{\tilde{E}}

\newcommand{\kapitel}{\chapter}

\newcommand{\proz}{^0\!/\!_{0}}
\newcommand{\prom}{^0\!/\!_{00}}

\renewcommand{\textfraction}{0.01}
\renewcommand{\floatpagefraction}{0.01}
\renewcommand{\topfraction}{0.99}
\renewcommand{\bottomfraction}{0.99}

%\renewcommand{\figurename}{\bf Fig.}
%\renewcommand{\tablename}{\bf Tab.}

\newcommand{\fig}[8]	%  8 parameters are required:
			% #1: FIGURE width: cm
			% #2: FIGURE hight: cm
			% #3: FIGURE caption width: cm
			% #4: figure position: hERE, tOP, bOTTOM
			% #5: figure label
			% #6: figure caption (TEXT ZUM BILD)
			% #7: entry for figure caption
			% #8: Bild File (Postscript)
{
\begin{figure}[#4]
   %\vspace{0.5cm}
   \centering{
   \epsfig{file=#8,width=#1cm}
   \parbox[t]{#3cm}
       {\vspace{-0.2cm} \caption[#7]{\label{#5} #6}}
   } % centering
   \vspace{-0.65cm}
\end{figure}
}

\newcommand{\figo}[9]	% 9 parameters are required:
			% #1: FIGURE width: cm
			% #2: FIGURE caption width: cm
			% #3: vertical offset preceeding picture: cm
			% #4: vertical offset following picture: cm
			% #5: figure position: hERE, tOP, bOTTOM
			% #6: figure label
			% #7: figure caption (TEXT ZUM BILD)
			% #8: entry for figure caption
			% #9: Bild File (Postscript)
{
\begin{figure}[#5]
   \vspace{#3cm}
   \centering{
   \epsfig{file=#9,width=#1cm}
   \parbox[t]{#2cm}
       {\caption[#8]{\label{#6} #7}}
   } % centering
   \vspace{#4cm}
\end{figure}
}

\newcommand{\abb}[4]	%  4 parameters are required:
			% #1: figure label
			% #2: figure caption (TEXT ZUM BILD)
			% #3: entry for figure caption
			% #4: Bild File (Postscript)
{
\begin{figure}[htb]
   %\vspace{0.5cm}
   \centering{
   \epsfig{file=#4,width=12cm}
   \parbox[t]{12cm}
       {\caption[#3]{\label{#1} #2}}
   } % centering
   \vspace{-0.5cm}
\end{figure}
}

% *** muessen unbedingt hinter \pagestyle{fancy} stehen, da dort defaults def.
%\renewcommand{\chaptermark}[1]{\markboth{\thechapter\ #1}{}}
%\renewcommand{\sectionmark}[1]{\markright{\thesection\ #1}}
%\renewcommand{\sectionmark}[1]{\markright{\thesection}}
%\renewcommand{\chaptermark}[1]{\markboth{#1}{}}
%\renewcommand{\sectionmark}[1]{\markright{#1}}

